\subsection{差距综合与研究定位}
\label{sec:gap-synthesis}

上述回顾确定了激光超声去噪中四个相互关联的研究差距。首先，常规去噪方法和通用深度学习模型优化SNR或MSE目标，而不对声学相关信号特征进行显式约束；因此，激进的噪声抑制通常会降低对缺陷检测至关重要的Lamb波到达时间、振幅比和波形形态 \cite{chen2019wavelet}。其次，仅在高平均数据上训练的深度学习模型在部署于低平均或含缺陷试件时表现出跨域泛化失败，其中信号统计与训练条件显著不同 \cite{liu2022deep, zhang2020deep}。第三，对广泛信号平均的依赖以实现可接受的SNR会施加线性时间惩罚，并有可能对敏感复合材料造成表面改性，然而减平均操作在文献中仍未得到充分探索。第四，尽管多域融合已显示出在导波SHM中改善健康指标提取的能力 \cite{perry2026health}，但在激光超声去噪中仍基本未被探索；现有深度学习方法主要在单一域中运行，无法联合捕获互补的时域和频谱信息。

这些差距在文献中由部分解决方案单独解决，但没有单一框架同时解决所有四个问题。本工作提出PMDF-Net，一个物理约束多域融合网络，通过声学物理损失项同时解决噪声抑制-特征保留的权衡，通过在完好和有缺陷试件上训练的缺陷敏感架构解决跨域泛化挑战，通过单次推理能力解决减平均约束，并通过在统一框架内融合时域、频域和时频表示解决多域信息差距。第\ref{sec:methodology}节详细描述了所提出的架构。

% word count: ~230
